\section{Evaluation Methodology}
\label{sec:method}

\begin{table}[t]
\centering
\caption{Schemes compared: 13 baselines, M2 and variants. All are implemented in Python; $\dagger$ marks schemes that are also in the Go prototype.}
\label{tab:baselines}
\footnotesize
\setlength{\tabcolsep}{3pt}
\begin{tabular}{@{}lp{6.6cm}@{}}
\toprule
ID & Scheme \\
\midrule
B1$^\dagger$ & HMAC-SHA256, key shared with all edges (EdgeAuth-like) \\
B2$^\dagger$, B3$^\dagger$, B4 & Ed25519 (Media CDN); ECDSA P-256 and RSA-2048 (CloudFront) \\
B5$^\dagger$, B6$^\dagger$, B7, B8 & per-token ML-DSA-44/65, SLH-DSA-128s/128f \\
B9$^\dagger$ & per-token hybrid Ed25519 + ML-DSA-44 \\
B10 & multi-MAC: one 16-B HMAC tag per verifier domain, $k\in\{1,3,10\}$ \\
B11$^\dagger$ & delegated Ed25519 key, certified per lifetime $L$ by ML-DSA-44 \\
B12$^\dagger$ & per-epoch HMAC key wrapped with ML-KEM-768~\cite{fips203} to each edge \\
B13 & MTL mode, one series per key~\cite{draft-mtl-mode} \\
\textbf{M2}$^\dagger$ & \textbf{stapled per-token ML-DSA-44} (Section~\ref{sec:design}) \\
\midrule
B5v$^\dagger$, B5c & B5 with an edge verify cache; B5 split over two cookies \\
O1$^\dagger$/O2/O3/O4$^\dagger$ & EMBER, $n=16$/$n=32$/SLH-DSA-128s root/hybrid root \\
\bottomrule
\end{tabular}
\end{table}

\textbf{Schemes.} Table~\ref{tab:baselines} lists the 13 baselines (B1--B13), M2 and the variants. Macaroons and DPoP are discussed only qualitatively (Section~\ref{sec:background}). In B13 the MTL series grows over $\vmax$. An MTL series restarted at every epoch is structurally the same as EMBER, so we do not list it separately.

\textbf{Testbed and implementations.} All experiments ran on one KVM guest with 4 vCPUs (Xeon at 2.1~GHz, AVX-512, SHA-NI) and 15~GB of RAM, running Linux~6.18. Primitive costs come from a native C benchmark against OpenSSL~3.5.4 (EVP interface, 5 fresh processes pinned to one core). We report the median of the run medians together with the [min--max] of the run medians. The Go prototype uses Go~1.24.7 and cloudflare/circl~1.6.1, and Python (pqcrypto~1.0.0) is a secondary source. We ran the native suite twice, separated by a VM reboot. Absolute values differ by 15--20\% (ML-DSA-44 verification took 100.9 vs.\ 122.3\us), while ratios are stable. All tables and all simulations use the second run, except the epoch-controller study and the simulator consistency check (Appendix~\ref{app:secondary}). Rerunning the parameter sweep with second-run costs left misses unchanged bit for bit and moved $D$ by at most 0.1~ms, so the controller conclusions, whose effects are all far below 10~ms, cannot change.

\textbf{Prototype.} The Go prototype consists of an edge (\texttt{net/http}; root cache with coalescing and a negative cache; optional verified-token cache; $V$ virtual edges), an issuer and root service, and an open-loop load generator. Client, edge and origin run in separate network namespaces joined by userspace link emulators (\texttt{netem} is unavailable): one-way delays of 10~ms and 37.5~ms, 1~Gbit/s, MTU 1500. \emph{P1} measures edge CPU cost. It is a closed loop over the edge namespace's loopback with the edge on one core, 64 keep-alive connections, $K=150$ requests per token and 5 runs. Its primary metric is edge CPU time per HTTP request. P1 measures CPU, not the network. \emph{P3} measures root dissemination over the emulated links with $V\in\{10,50\}$ virtual edges, $\Delta\in\{0.1,1,10\}$~s and $\lambda\in\{20,200\}$/s, in pull and staple-first modes: 12 configurations, each with a 60~s window and 5 runs. The prototype has no M2 or staple-all dissemination run and no pooled-EMBER run. \emph{P2} measures latency under load (Appendix~\ref{app:secondary}). TLS is not included.

\textbf{Wire bytes and real clients.} Request-header bytes per session (1 manifest and 150 segment requests) are encoded with Go's HPACK at 4096-B and 64~KiB tables. To observe real encoders, a raw-frame HTTP/2 server advertises a table size and logs each request's header-block length while a client fetches 150 objects with a \num{3311}-B cookie value. The clients are headless Chrome~153 (3 runs), curl~8.5.0/nghttp2~1.59.0 and the Go x/net client. We did not test Firefox or Safari.

\textbf{Fleet simulator.} The simulator is an epoch-level Monte Carlo engine driven by the native costs. Its default fleet has $E=500$ edges with Zipf(0.8) traffic shares and 50 shields, and $T=1000$ tenants with Zipf(1.0) rates and aggregate rate $\bar\lambda\in\{200,2000,\num{20000}\}$/s. Sessions have $K$ geometric with mean 150, 4~s apart, and each segment is re-routed to another edge with probability $p_{\text{move}}=0.0005$ (about 7\% of sessions move once). RTTs are lognormal: edge--shield 15~ms, shield--root 60~ms, token receipt to first request 120~ms. All of these are assumptions; the sensitive ones are swept (Appendix~\ref{app:secondary}). The pooling policies are per-tenant trees, 8 shards, 10 or 50 regional pools, and one pool per issuer key. Each is run with a fixed $\Delta$, work-conserving closing (WC, close whenever the signer is idle) or M2. CIs are Student-$t$ 95\% intervals over 5 seeds with common random numbers. Unless stated, their half-widths are below 1\% of the mean.

\textbf{Hypotheses.} We fixed hypotheses H1--H9 in the design document before running any experiment (repository commit \texttt{3148c65}, 2026-11-05 18:07 UTC). Appendix~\ref{app:hyp} reports every outcome.
